\section{Interpretation, limitations, and recommendations}
\label{sec:discussion}

The principal outcome is a reviewable mathematical archive produced
under a broad, fixed-duration instruction. It contains explicit
constructions, arbitrary-size arguments, finite certificates, and
records of reclassification. The experiment demonstrates that this
particular configured agent could sustain such a workflow across many
continuations. It does not by itself establish a general success rate
for open problems, a comparison with professional mathematicians, or
the causal contribution of any single component.

\subsection{What the evidence can support}

The strongest local evidence is a short argument whose hypotheses and
conclusion can be inspected directly. The examples in
Section~\ref{sec:examples} are included for that reason. Finite
computations can add exact evidence, especially when their output is a
certificate checked in a different representation. However, two
implementations written by the same agent may share a mistaken
mathematical interpretation. A passing checker establishes only what
its implemented predicate and coverage justify. The corrected
\kn{21.52} test supplies a concrete instance of this limitation.

Likewise, an explicit reduction from the printed question to an
imported theorem may be more informative than a long new calculation.
Its obligations are to cite the correct version, check every
hypothesis, and account for conventions. It should be described as a
deduction from prior work when that is what it is. The \kn{17.34}
episode shows that further literature review can change the novelty
classification without invalidating the mathematics. Such changes
belong in the results of a research experiment.

The archive also shows practical value in persistent written state.
Proofs, process records, and explicit boundaries survived between
turns; unfinished checks remained identifiable; later source findings
could be propagated into a consolidated ledger. This is an observation
about how the run proceeded. Without an ablation, it cannot quantify
how much persistence, Git, browsing, or GAP improved the outcome.
The repeated user-facing references to GAP in continuation messages
also suggest that the conversation's local cues affected the wording
of updates; they should not be read as independent installation events.

\subsection{Limits of the case study}

Selection was highly nonuniform. The agent had access to the entire
public Notebook and chose entries opportunistically. It could recognize
a familiar construction, find an implication in a recent paper, or
abandon an expensive line. The 46-entry list is therefore an outcome
inventory, not 46 successes in a prespecified sample of 1,308 trials.
The number of meaningfully attempted targets cannot be recovered simply
by counting indexed entries or occurrences of problem numbers in logs.

Training-data overlap is unknown. Many questions and ingredients were
public long before the run. Even a proof that does not quote a source
may use knowledge learned from it. Conversely, an elementary argument
can be independently rediscovered after its ingredients become familiar.
Neither possibility can be decided from the wording of an answer.
The experiment had browsing access, so retrieval and mathematical
reasoning were deliberately interwoven. We make no claim to have
isolated either capability.

The full JSONL improves behavioral reconstruction through dated tool
observations, request identifiers, compaction records and corroborated
model settings. It remains a client-side record: remote inference
hardware, actual peak local resource consumption, provider invoices,
and a complete historical sampling configuration are unavailable.
The cost reconstruction in Section~\ref{sec:methods} makes its rates
and counter scope explicit; it does not measure actual billing or
energy use. Human oversight outside the recorded interaction is also
not observable from the trace.

The initial manuscript review was performed by the same agent that
produced the research. The deadline archive contains zero outside
reviews. A later external review and follow-up checked arguments and
finite inputs, identified scope and packaging issues, and withdrew
several initial objections. We preserve the correspondence in
Appendices~\ref{app:review-original}--\ref{app:reply-final}, with a
change record in Appendix~\ref{app:review-changes}. The reviewer's
scripts and logs are separately versioned; their reported outcomes
are distinguished from checks rerun during this revision.
The revised 14.72 construction makes the fixed divisor globally
principal, and the portable bundle now includes the missing small
certificate. These are later changes, not deadline achievements.
Neither a preserved PASS message nor a reviewer's acceptance
recommendation substitutes for specialist assessment or establishes
priority. Codex \texttt{gpt-6-astra} is
credited as the principal author for research execution, proof and code
drafting, manuscript preparation, and internal review. Michael Kinyon
and Josef Urban conceived the prompt and experiment and oversaw the
work. These contributions do not constitute outside mathematical review.

\subsection{A more informative follow-up experiment}

A follow-up should retain the broad research setting while improving
the measurements that this run could not provide. A practical protocol
would have the following components.

\begin{enumerate}
\item \textbf{Define the unit of evaluation in advance.}
Use a dated set of problems with explicit subparts and a documented
editorial status. Record every substantive attempt, including abandoned
ones, with the reason for stopping. A randomized target subset or
matched target sets would support comparisons that an opportunistic
portfolio cannot.

\item \textbf{Separate discovery from assessment.}
Maintain different fields for statement coverage, proof validity,
computational verification, prior-art status, and outside review.
Freeze proposed answers before specialist assessment, and preserve
later corrections as new versions. Reviewers should receive the
original question and the mathematical argument before seeing aggregate
success claims.

\item \textbf{Record the experimental configuration.}
Save model identifiers, reasoning settings, harness and package
versions, continuation instructions, tool permissions, and retrieval
dates. Capture raw event logs with timestamps and request identifiers,
along with local CPU time and peak memory. Retain response-level usage, including compaction, and separate the
research interval from subsequent writing. Keep usage estimates and
billing records separate until reconciled.

\item \textbf{Treat verification as a research task.}
For a search, specify the exact predicate and its relation to the
question. Include a small direct implementation that checks examples
without reusing the search representation. For a large certificate,
make coverage and the checking rules explicit, retain malformed-input
and negative controls, and distinguish successful generation from
successful verification.

\item \textbf{Use formalization selectively.}
Small algebraic identities, finite certificates, and critical reduction
lemmas are plausible first targets. Formalizing the entire archive at
once would create a large translation project. Any formal account must
also justify that its statement matches the Notebook, including
quantifiers and conventions; a kernel-checked proof of a different
statement would leave the original task unanswered.

\item \textbf{Measure the review burden.}
Record expert time, corrections, changes in status, and which claims
depend on the same imported theorem. Report a short, fully reviewed
core alongside the larger exploratory archive. This would make it
possible to compare output volume with the work needed to establish
reliable mathematical contributions.
\end{enumerate}

The present archive is suited to such a second phase. Its most useful
next product would be a set of independently assessed individual
results, with corrected proofs and precise attribution, followed by
formal certificates where they materially reduce the remaining trust.
The current paper provides the experimental context and mathematical
entry points for that work.
